\documentclass[11pt,a4paper]{article}
\usepackage[margin=2.2cm]{geometry}
\usepackage[T1]{fontenc}
\usepackage[utf8]{inputenc}
\usepackage{lmodern}
\usepackage{booktabs}
\usepackage{array}
\usepackage{xcolor}
\usepackage{enumitem}
\usepackage{hyperref}
\usepackage{microtype}
\hypersetup{colorlinks=true,linkcolor=blue!60!black,citecolor=blue!60!black,urlcolor=blue!60!black,breaklinks=true}
\setlist{nosep,leftmargin=*}
\newcommand{\model}{\textsc{TomaNet}}
\newcommand{\modelcls}{\textsc{TomaNet-C}}
\newcommand{\modeldet}{\textsc{TomaNet-D}}

\title{\textbf{Our Approach: \model{}}\\[4pt]
\large A Plain-English Guide to the Model We Built}
\author{Tomato disease detection project}
\date{\today}

\begin{document}
\maketitle
\tableofcontents
\newpage

\section{What This Project Is}
Tomato plants get sick from things like bacterial spot, blight, and viruses. A farmer or
a phone app can either:
\begin{itemize}
  \item look at a photo of \emph{one leaf} and say \emph{what disease this is} -- we call
    this \textbf{classification}, or
  \item look at a photo of a \emph{whole plant} and \textbf{draw a box} around every
    diseased spot it can find, even several at once, on several leaves -- we call this
    \textbf{detection}.
\end{itemize}
We built one small, efficient model design that does both jobs, and the data pipeline
needed to train and check it honestly.

\section{The Model: \modelcls{} and \modeldet{}}
Both are built from the same reusable building block, called \textsc{TomaBlock}. Think of
it like a Lego brick: the exact same brick is used to build the ``trunk'' -- the part
that looks at the raw image and figures out what's in it -- for \emph{both} tasks.

\begin{itemize}
  \item \textbf{\modelcls{}} (classification) = trunk $\to$ one small extra layer that
    turns the trunk's output into a single disease name.
  \item \textbf{\modeldet{}} (detection) = the exact same trunk $\to$ a ``neck'' that
    combines information at different zoom levels $\to$ a head that draws boxes.
\end{itemize}

Because the trunk is identical -- not just similar, literally the same list of layers,
checked automatically by an automated test -- a trunk trained on one task can jump-start
training on the other. This is not a marketing claim; it is a structural fact about the
code.

Both come in three sizes, so the same design scales from a phone to a more capable
server. These are the real, measured numbers (not estimates) from actually running our
model-building script:

\begin{center}
\begin{tabular}{@{}lccc@{}}
\toprule
 & \textbf{n (smallest)} & \textbf{s} & \textbf{m (largest)} \\
\midrule
\modelcls{} (classification) & 1.32\,M params & 4.47\,M params & 14.44\,M params \\
\modeldet{} (detection) & 2.92\,M params & 10.64\,M params & 30.79\,M params \\
\bottomrule
\end{tabular}
\end{center}
\vspace{2mm}
\noindent (``M'' = million. For comparison, popular phone-sized detectors are usually in
the same 1--5\,M range.)

\section{What Makes Our Design Different}
We are not claiming to have invented brand-new ideas from scratch. We took three good,
published ideas and combined them in a way nobody had applied to tomato disease work
before, and we picked each one for a specific, explainable reason.

\begin{enumerate}
  \item \textbf{We only convolve a quarter of the channels at a time (``PConv'').}
    A normal convolution touches every channel of every layer, which is slow mainly
    because of memory movement, not the actual math -- and memory movement is the real
    speed bottleneck on small devices like a phone. This idea comes from a method called
    FasterNet \cite{fasternet}. We use it so the model stays fast and light.
  \item \textbf{We add a positional attention step (``Coordinate Attention'').}
    Disease spots on a leaf are small and appear in specific places, not spread evenly
    across the whole leaf. A common shortcut used in many lightweight models
    (channel-only attention) throws away exactly that positional information.
    Coordinate Attention \cite{coordatt} keeps track of \emph{where} on the leaf
    something unusual is happening, which matters for small, localised problems like
    lesions.
  \item \textbf{We train with a richer version of the block, then simplify it before
    deployment (``re-parameterisation'').}
    During training, our block briefly has three small parallel paths, which helps it
    learn better. Right before the model is used for real, those three paths are
    mathematically folded back into one single, simple path -- so training gets the
    benefit, but the final deployed model pays no extra cost. This RepVGG-style trick
    \cite{repvgg} is checked in our own tests to be numerically exact (the folded and
    unfolded versions agree to about five decimal places).
  \item \textbf{One trunk, two jobs.}
    As described above, the same trunk serves both classification and detection,
    byte-for-byte identical in the model definition. This is the part of the design that
    is genuinely uncommon in tomato disease research specifically -- most published work
    builds a completely separate model for each task.
\end{enumerate}

\noindent We are honest about scale: our detection model is not smaller than every
competing lightweight detector on the market. The point we are actually trying to prove
is that one small, shared design can do both jobs well and transfer knowledge between
them -- and that claim has to be measured with real experiments, not assumed.

\section{Datasets We Actually Use}
Many public tomato datasets exist. Only the ones below are wired into working code right
now -- meaning there is a script that downloads them, cleans them, and prepares them for
training. Nothing else is included here; this list is deliberately limited to what
actually runs today.

\subsection{For Classification (``what disease is this leaf?'')}
\begin{itemize}
  \item \textbf{PlantVillage} -- 18{,}160 tomato leaf photos, 10 disease/healthy classes.
    Taken in a lab on a plain background, not in a real field.
  \item \textbf{Taiwan tomato dataset} -- 622 real field photos, 6 classes. Small, but
    real-world, so it is our first sanity check for any new model.
  \item \textbf{PlantDoc} -- around 700 tomato photos scraped from the web, 9 classes.
  \item \textbf{CCMT} -- field photos from Ghana, 5 classes.
  \item \textbf{PlantVillage (Zenodo re-split)} -- a smaller, pre-split copy of
    PlantVillage, used as a quick stand-in when a full run isn't needed.
\end{itemize}

\subsection{For Detection (``draw a box around every disease spot on this plant'')}
\begin{itemize}
  \item \textbf{Tomato-Village (detection variant)} -- field photos from Rajasthan,
    India. This is currently our \emph{only} working detection dataset. After preparing
    it, we have 14{,}259 usable images with 6 disease classes that actually have boxes
    drawn on them. (Two of the eight label names in the source data turned out to have
    no boxes anywhere in the files -- a real property of the raw data, confirmed by
    checking every label file directly, not a mistake in our code.)
\end{itemize}

\noindent No other public dataset has a working detection adapter yet. We checked one
strong candidate, PlantDoc, and found that its public copy on GitHub does not actually
contain any box annotation files, despite being listed elsewhere as box-annotated -- so
it is classification-only in practice, at least from that source.

\subsection{An Important Data-Quality Fix We Made}
The Tomato-Village detection data ships with its own train/test split already done by
its authors -- but we found that split leaks: many photos in their ``test'' portion are
slightly-edited duplicates of photos in their ``training'' portion (rotated, recoloured,
etc.), which would make any accuracy number misleadingly high. We wrote our own,
corrected split that keeps every edited copy of the same original photo together on one
side, so the test numbers we report are trustworthy.

\section{Honest Limits -- What We Have Not Done Yet}
What exists today is: a working, tested data pipeline for one classification job and one
detection job, a shared model architecture proven to build correctly at all six sizes,
and one real (if short) training run on each task confirming the whole pipeline works end
to end. Larger training runs, comparisons against other models, and the deeper research
questions (cross-domain generalisation, explainability, unknown-disease detection) are
the next steps -- see the companion document, \emph{Research Context}, for what the
literature says about those questions and how they shape our plan.

\begin{thebibliography}{9}
\footnotesize
\bibitem{fasternet} J.~Chen et al., ``Run, don't walk: partial convolution (FasterNet),'' \emph{CVPR}, 2023. arXiv:2303.03667
\bibitem{coordatt} Q.~Hou, D.~Zhou, J.~Feng, ``Coordinate attention for efficient mobile network design,'' \emph{CVPR}, 2021. arXiv:2103.02907
\bibitem{repvgg} X.~Ding et al., ``RepVGG: making VGG-style ConvNets great again,'' \emph{CVPR}, 2021. arXiv:2101.03697
\end{thebibliography}

\end{document}
